\chapter{Back Pain}

\section*{Definition}
Back pain is discomfort in the back, most often in the thoracic or lower-back region. It may arise from muscles, joints, discs, nerves, bones, or internal organs. This chapter focuses mainly on back pain in adults; neck pain is discussed separately.

\section*{When to See a Doctor}

\begin{itemize}
 \item Seek emergency care for new loss of bladder or bowel control, inability to urinate, numbness around the genitals or anus, or rapidly progressive leg weakness; these can indicate cauda equina or spinal cord compression.
 \item Seek prompt medical assessment for fever or feeling very unwell, major trauma, unexplained weight loss, a history of cancer, significant immune suppression, osteoporosis, or long-term steroid use.
 \item Arrange medical assessment for persistent, recurrent, unexplained, or night-waking pain, or pain associated with abdominal, urinary, or other systemic symptoms.
\end{itemize}

\section*{How It Develops}
Back pain may be nonspecific, with no single structural cause identified, or it may result from a specific spinal or nonspinal condition. It can be acute, subacute, or chronic. Pain that varies with movement is often mechanical, whereas prolonged morning stiffness, night pain, and improvement with activity may suggest inflammation and need clinical assessment.

\section*{Causes}
\begin{itemize}
 \item Muscle strain
 \item Herniated disc
 \item Degenerative disc disease
 \item Osteoarthritis
 \item Spinal canal narrowing (spinal stenosis)
 \item Osteoporotic fracture
 \item Infection of a vertebra or disc
 \item Inflammatory arthritis, including axial spondyloarthritis
 \item Cancer affecting the spine
 \item Referred pain from the kidneys, abdominal organs, or blood vessels
\end{itemize}

\section*{Related Symptoms}
\begin{itemize}
 \item Stiffness
 \item Limited mobility
 \item Radiating leg pain
 \item Numbness, tingling, or weakness
 \item Muscle spasms
 \item Reduced flexibility
 \item Pain or stiffness that is worse after rest or in the morning
\end{itemize}

\section*{Medical Evaluation}
\begin{itemize}
 \item History and examination assess neurologic function, gait, range of motion, systemic symptoms, trauma, and red flags.
 \item Imaging is not routinely needed at first for uncomplicated back pain. It is considered when red flags, major trauma, osteoporosis, or a significant or progressive neurologic deficit raises concern for a specific cause, and when the result is likely to change management.
 \item MRI is commonly used when serious nerve compression, infection, cancer, or persistent/progressive neurologic symptoms are suspected; radiographs or CT may be used when a fracture or bony problem is suspected.
 \item Blood tests are considered when infection, inflammatory disease, malignancy, or another systemic cause is suspected.
\end{itemize}

\section*{Treatment Options}
\begin{itemize}
 \item Most nonspecific back pain is managed with education, reassurance, continued activity, exercise or physical therapy, and selected analgesics.
 \item An oral NSAID may be considered when safe, using the lowest effective dose for the shortest time; clinicians assess gastrointestinal, kidney, liver, heart, pregnancy, and medication-related risks and may prescribe gastroprotection.
 \item Radiculopathy, fracture, infection, inflammatory disease, or malignancy requires cause-specific treatment and sometimes specialist referral.
 \item Procedures or surgery are reserved for selected conditions, such as persistent disabling sciatica with matching imaging findings or emergency nerve compression.
\end{itemize}

\section*{Self-Care and Home Management}
\begin{itemize}
 \item Stay gently active and avoid prolonged bed rest; temporarily modify activities that clearly worsen pain.
 \item Heat or ice may provide short-term relief. Do not apply either directly to the skin or while asleep.
 \item Ask a clinician or pharmacist before using an NSAID if you have kidney disease, a history of ulcer or gastrointestinal bleeding, heart disease, are taking blood thinners, or are pregnant. Do not combine NSAIDs unless advised.
 \item Do not use paracetamol/acetaminophen alone as the main treatment for low back pain, as it may not be effective; ask a clinician about alternatives if an NSAID is unsuitable.
 \item Seek urgent care for new bladder or bowel dysfunction, inability to urinate, saddle numbness, rapidly progressive leg weakness, fever with severe back pain, or major trauma.
\end{itemize}

\section*{References}
\begin{thebibliography}{99}
\bibitem{back_pain:ref1} National Institute for Health and Care Excellence. ``Low back pain and sciatica in over 16s: assessment and management (NG59).'' Updated 2020. https://www.nice.org.uk/guidance/ng59.
\bibitem{back_pain:ref2} Deyo RA, Weinstein JN. ``Low back pain.'' \textit{N Engl J Med}. 2001;344(5):363-370.
\bibitem{back_pain:ref3} Patrick N, Emanski E, et al. ``Acute and chronic low back pain.'' \textit{Med Clin North Am}. 2016;100(2):351-366.
\bibitem{back_pain:ref4} Oliveira CB, Maher CG, et al. ``Clinical practice guidelines for the management of non-specific low back pain.'' \textit{Eur Spine J}. 2018;27(11):2791-2803.
\bibitem{back_pain:ref5} American College of Radiology. ``ACR Appropriateness Criteria: Low Back Pain.'' Revised 2021. https://acsearch.acr.org/docs/69483/Narrative/.
\bibitem{back_pain:ref6} World Health Organization. ``WHO guideline for non-surgical management of chronic primary low back pain in adults in primary and community care settings.'' Geneva: WHO; 2023.
\end{thebibliography}
